\section{An economy with interior saving}\label{sec:interior}

The corner economy keeps the cross-section on two asset levels, so its history is short. Here the
cross-section carries its whole history, saving is interior, and $(z,K)$ is still sufficient,
because capital encodes the history. The route is the order route of Section~\ref{sec:aggregate},
with an economy in which its hypotheses can be checked.

\subsection{Separation from ordered cash}

Three conditions make the two hard hypotheses of Theorem~\ref{thm:order} free.
\begin{itemize}[leftmargin=*,itemsep=2pt]
\item \textbf{Prices depend on the shock only}, as with a linear technology $Y=A(z)K+B(z)L$, whose
  marginal products are exact. Then the household's problem does not depend on the path of
  capital, so saving rises (weakly) with capital, and sequential equilibrium is unique: the
  distribution path does not depend on prices.
\item \textbf{Shocks are i.i.d.}: next period's aggregate shock and employment are drawn
  independently of today's. Then the value of next period's assets is the same function at the
  root of every history, so today's saving is one nondecreasing function of cash on hand (a
  Topkis argument across histories).
\item \textbf{Cash is ordered by the shock}: after a worse shock every household has less cash than
  every household after a better one.
\end{itemize}
Under these, every household saves less after the worse shock, so the next cross-sections are
ordered for any two current ones. This is the separation condition of the order route, directly,
without the extreme histories.

\begin{theorem}\label{thm:orderedcash}
If prices depend on the shock only, shocks are i.i.d.\ and cash on hand is ordered by the shock,
then in every sequence-state equilibrium $(z,K)$ is a sufficient state.
\end{theorem}
\formal{kSufficient\_of\_orderedCash, sepOnHistories\_of\_orderedCash, policy\_le\_policy\_of\_cont,
cont\_root\_eq, uniqueSequentialEquilibrium\_of\_exogenous}

\subsection{Interior saving}

Interiority is proved by two corner deviations. Extra assets $\delta$ are worth at least
$R\delta/(\text{cash}+R\delta)$ next period: keep the saving of the poorer state and consume the
extra return, and use the tangent line of the logarithm. The last $\delta$ below $b$ is worth at
most $R\delta/(\text{cash}(b)-\abar-R\delta)$: the poorer state copies the richer one's saving. So
saving is positive at cash $x$ if
\[
  \frac{\delta}{x-\delta} < \beta\,\E\!\left[\frac{R'\delta}{w'e'+R'\delta}\right],
\]
and below the cap if
\[
  \frac{\delta}{x-\abar+\delta} > \beta\,\E\!\left[\frac{R'\delta}{(R'-1)\abar+w'e'-R'\delta}\right].
\]
\formal{value\_sub\_ge, value\_sub\_le, policy\_pos\_of, policy\_lt\_cap\_of}

\subsection{The economy}

\begin{center}
\begin{tabular}{ll}
\toprule
Aggregate shocks & two, i.i.d.\ with probability $1/2$ \\
Efficiency & $1$ or $3/4$, i.i.d.\ with probability $1/2$ \\
Technology & linear: gross return $6/5$; wage $1$ (good), $1/3$ (bad) \\
Preferences & logarithmic, $\beta=1/2$; asset cap $\abar=25/72$; floor $1/100$ \\
\bottomrule
\end{tabular}
\end{center}

\noindent
After a bad shock cash is at most $\tfrac65\cdot\tfrac{25}{72}+\tfrac13=\tfrac34$; after a good
one it is at least $\tfrac34$. With $\delta=1/3000$, for low-efficiency households after a good
shock (cash in $[3/4,\,7/6]$) the first condition reads $1/2249 < \tfrac18(1/2501+1/1876+3/2503+1/626)$
and the second $3/7378 > \tfrac18\cdot18\,(1/48107+1/36857+1/18107+1/14357)$, both in exact
rational arithmetic.

\begin{theorem}\label{thm:interior}
This economy has a sequence-state equilibrium. In every one, $(z,K)$ is a sufficient state, and
after every good shock at least half of the households --- all those with low efficiency, whatever
their assets --- save strictly between $0$ and $\abar$.
\end{theorem}
\formal{InteriorWitness.exists\_interiorWitness, InteriorWitness.interiorWitness,
InteriorWitness.saving\_interior, InteriorWitness.orderedCash}

\noindent
The costs are plain. Prices do not respond to capital, so the equilibrium feedback that makes
Krusell--Smith hard is absent; what remains is the aggregation question itself, whether aggregate
capital pins down a cross-section that keeps its whole history, and here it does. Shocks are
i.i.d.\ rather than persistent. And the household is impatient: with these crude interiority
bounds the windows are tiny for $\beta\ge 0.8$, so the period should be read as long. We do not
prove that interior saving is strictly increasing in assets, which would make the cross-section
infinitely supported.
